\title{XTab: Cross-table Pretraining for Tabular Transformers}

\begin{abstract}
The success of self-supervised learning in computer vision and natural language processing has motivated pretraining methods on tabular data. However, most existing tabular self-supervised learning models fail to leverage information across multiple data tables and cannot generalize to new tables. In this work, we introduce XTab, a framework for cross-table pretraining of tabular transformers on datasets from various domains. We address the challenge of inconsistent column types and quantities among tables by utilizing independent featurizers and using federated learning to pretrain the shared component. Tested on 84 tabular prediction tasks from the OpenML-AutoML Benchmark (AMLB), we show that (1) XTab consistently boosts the generalizability, learning speed, and performance of multiple tabular transformers, (2) by pretraining FT-Transformer via XTab, we achieve superior performance than other state-of-the-art tabular deep learning models on various tasks such as regression, binary, and multiclass classification.
\end{abstract}

\section{Introduction} 
\label{introduction}
With the increasing number of datasets represented as tables with rows and columns, tabular machine learning makes the foundation of many real-world applications. While deep learning has achieved tremendous success in the fields of computer vision (CV) \citep{he2022masked, liu2021swin} and natural language processing (NLP) \citep{devlin2018bert, vaswani2017attention}, tabular deep learning models are not used as commonly as tree-based models~\citep{grinsztajn2022tree, gijsbers2022amlb}. 

The primary challenge of tabular deep learning is the diversity of tabular tasks. Unlike text, which can be standardized as a sequence of tokens, tables are highly data-specific. Tabular data can vary in the number and types of columns. This makes it difficult for tabular deep learning models to transfer the knowledge learned from one table to another, leading to poor generalization abilities. Therefore, self-supervised learning for tabular data~\citep{he2022masked, devlin2018bert}, particularly one that is able to bootstrap the learning on new tables, is still an open problem.

There is an ongoing effort in migrating self-supervised pretraining techniques from CV~\citep{chen2020simple} and NLP~\citep{devlin2018bert} to tabular tasks.

With self-supervised pretraining, tabular deep models have demonstrated improved performance~\citep{ucar2021subtab, bahri2021scarf, majmundar2022met}. However, existing methods generally pretrain the tabular model on data from the same domain as the downstream task. As a result, the data-specific models cannot generalize to new tables.

Another direction of deep tabular learning aims to leverage Transformers, which drives the recent progress in NLP~\citep{vaswani2017attention} and CV~\citep{dosovitskiy2020image} for tabular tasks. Inspired by the success of the attention mechanism, Transformers were adapted to tabular data~\citep{gorishniy2021revisiting, somepalli2021saint, wu2021fastformer, wang2022transtab} and demonstrated strong performance~\citep{grinsztajn2022tree}. The core idea of tabular transformers is to consider the table columns as tokens, similar to words in a sentence. Therefore, tabular transformers can process tables with variable numbers of columns, thus making transferable learning \citep{wang2022transtab} feasible.

In this paper, we present \textit{XTab}, a general framework for \textit{cross-table pretraining of tabular transformers}.

To resolve the issue that tables may vary in the number and types of columns, XTab decomposed the tabular transformers to two components: data-specific featurization and projection layers that capture the characteristics of each table, and a cross-table-shared block that stores the common knowledge. On a diverse collection of data tables, XTab trains these data-specific blocks and the shared block jointly via federated learning~\cite{collins2022fedavg}. Once pretrained, XTab can bootstrap the learning process on a new table by initializing the shared block with pretrained weights.  To verify our design, we conducted extensive experiments on AutoML Benchmark (AMLB)~\cite{gijsbers2022amlb}. Our results show that transformers pretrained and initialized with XTab consistently outperform transformers with random initialization. By pretraining FT-Transformer~\citep{gorishniy2021revisiting} with XTab, we outperform the state-of-the-art tabular deep learning models.

The contributions of the paper are summarized as follows:

\begin{itemize} 
\itemsep-0.3em 
\item XTab offers a framework to account for cross-table variations and enable cross-table knowledge transfer. 
\item Given the large diversity of tabular datasets, we propose to pretrain on tabular datasets with federated learning. This allows us to perform distributed pretraining across a large collection of tables. 
\item To the best of our knowledge, we are the first to show that cross-table pretraining can boost the learning speed and performance on new tables. This is different from table understanding tasks~\cite{yin2020tabert},  the focus of which is to extract the semantical information from tables.

\end{itemize}

\section{Related work} \label{related_work} 
\paragraph{Tabular self-supervised learning.} Inspired by the success of pretraining in CV and NLP, previous papers studied tabular self-supervised learning \citep{yoon2020vime, ucar2021subtab, somepalli2021saint, bahri2021scarf, majmundar2022met, rubachev2022revisiting, wang2022transtab}. Among those works, \citet{yoon2020vime, ucar2021subtab} proposed an auto-encoder framework with a pretext task to reconstruct the missing part of a table. \citet{bahri2021scarf} used contrastive learning as the pretraining objective and extended the SimCLR framework~\citep{chen2020simple} to tabular tasks. \citet{rubachev2022revisiting, wang2022transtab} further incorporated the label columns of tabular tasks in pretraining and proposed ``target-aware" objectives leading to higher performance. As existing approaches only pretrain on one \citep{bahri2021scarf,ucar2021subtab} or a few relevant tables \citep{wang2022transtab}, the pretrained tabular model lacks generalizability. XTab alleviates this issue by pretraining on a large number of tables. 

\paragraph{Tabular transformers.} Transformer models are gaining popularity in the realm of deep learning for tabular data. For example, FT-Transformer has demonstrated superior performance on tabular classification/regression tasks~\citep{gorishniy2021revisiting}. Saint introduces the row-wise attention and captures the inter-sample interactions using transformer~\citep{somepalli2021saint}. Fastformer proposes to use additive attention on tabular tasks, which is a lightweight attention mechanism with linear complexity to the length of input sequences~\citep{wu2021fastformer}. TransTab features transfer learning in tabular tasks using transformers~\citep{wang2022transtab} and also supports the cross-table transfer. Our approach is different from TransTab in that TransTab has limited ability in generalizing to tables from new domains, while XTab is able to generalize to new domains.

\paragraph{Cross-table transfer learning.} Pretrained vision and text models can be adapted to a wide range of tasks~\cite{bommasani2021opportunities}. One reason is that the sentences and images share general representations across various tasks. As for tabular learning, one may question if there is shared knowledge across tables as two different tables can have totally different numbers of columns and the associated semantic meanings. We argue that different tables share a similar prior given the recent success of zero-shot hyperparameter optimization (HPO) in AutoML~\citep{winkelmolen2020practical}, which learns a general hyperparameter configuration applicable to a wide range of tabular tasks. Unlike pretrained models in NLP~\citep{devlin2018bert}, XTab does not attempt to learn a universal tokenizer for all tables, as the meaning and context of each table varies. Instead, we aim to learn a weight initialization that is generalizable to various downstream tasks. Concurrent to our work, tabular prior-data fitted networks~(TabPFN)~\citep{hollmann2022tabpfn} learns a prior model on synthetic tabular data and demonstrated promising results on small numerical tabular classification tasks with $\leq1000$ samples. Different from TabPFN, the inference complexity of XTab is irrelevant to the number of training samples. Thus, XTab also works for large tables.

\section{Methods} 
Previous works have proposed various pretraining methods for tabular learning~\citep{bahri2021scarf, ucar2021subtab, rubachev2022revisiting, somepalli2021saint}. However, existing pretrained models are still domain-specific since they were pretrained on the training set of each individual tabular prediction task. As a result, existing pretrained models lack generalizability and fail to cover downstream tasks on other types of tables. Here, we propose XTab to pretrain transformer models using the information from multiple tables. With cross-table pretraining, XTab aims to learn the shareable knowledge that can boost the performance for various downstream regression and classification tasks.

\subsection{Model structure} The model structure of XTab is described in Figure \ref{model_structure}. During the pretraining phase, we sample mini-batches of rows from different tables (one batch per table). The featurizers are data-specific and convert each column of the table to a token embedding. An additional $\textsf{[CLS]}$ token is appended during this step for supervised prediction or contrastive self-supervised pretraining \citep{wang2022transtab}. A transformer-based backbone is shared across all tabular datasets to process token embeddings with variable sequence lengths. The output of the shared backbone is further processed by projection heads to (1) reconstruct the original table from a corrupted view; (2) identify the positive/negative pairs of samples as in contrastive learning; or (3) predict the values in the label column predefined by each table. The projection heads are not shared across tables since they are specific to each dataset and the pretraining objectives. Among all  pretraining losses, reconstruction loss and contrastive loss do not require information from the label column, whereas supervised losses use the groundtruth data in the label columns of each table. Using groundtruth information during the pretraining phase is referred to as ``target-aware pretraining"~\citep{rubachev2022revisiting, wang2022transtab} or ``pre-finetuning" \citep{aghajanyan2021muppet} in previous works.

A key challenge in cross-table pretraining lies in the variations of input tables. Previous works on transferable tabular learning either require tables to come from similar domains \citep{levin2022transfer} or use additional information (e.g., column names) to identify the shared knowledge across tables. XTab is designed to be applicable to previously unseen tables with no assumption on the domain or column name format. To this end, XTab contains  model blocks that carry the data-specific information (green blocks in Figure~\ref{model_structure}), as well as the shared backbone that stores the common knowledge (grey blocks in Figure~\ref{model_structure}). Once pretrained, only a shared backbone is kept for all downstream tasks. For each downstream task, featurizers and projection heads are randomly initialized and the entire model is finetuned on the downstream training data until a stopping criterion is met.

\subsubsection{Featurizers} \label{Featurizers} 
The featurizers convert a sample to feature embeddings $E \in \mathbb{R}^{c \times d}$. Here, $c$ denotes the number of columns and $d$ is the embedding dimension. Each row of a table is considered as an input sample, and each column is a token. The embedding of $\textsf{[CLS]}$ token is appended to the feature embedding for prediction $\text{stack}[E, \textsf{[CLS]}] \in \mathbb{R}^{c+1 \times d}$. In this work, we limit our discussion to tables with numerical and categorical columns. Text cells are treated as categorical attributes. Our tokenizer is similar to \citet{gorishniy2021revisiting}. For numerical features, we multiply the numerical value $x_{k}$ at the $k$-th column with a trainable vector $W_k \in \mathbb{R}^{d}$ and add a bias term $b_k$. For categorical columns, XTab learns an embedding matrix $\in \mathbb{R}^{N_{cat} \times d}$ as a lookup table, where $N_{cat}$ is the total number of categories of the dataset. During the forward pass, we retrieve the categorical feature embeddings from the embedding matrix.

XTab allows tables to have different numbers of columns and arbitrary column types. Featurizers are data-specific to handle various types and numbers of columns in the input.

\subsubsection{Backbones} \label{method_backbone} 
As the shared component across multiple pretraining datasets, transformers can handle input sequences with variable lengths. Therefore, it is possible to pretrain a tabular transformer that can be applied to all tabular datasets. Compared with other deep learning architectures like multi-layer perceptron (MLP), transformers are  favorable for cross-table knowledge transfer since they can handle variable input sequences~\citep{wang2022transtab}. As long as the backbone can process input sequences of variable lengths, XTab is flexible on the exact implementation. In this work, we present three  backbone variants: 

\textbf{FT-Transformer:} Feature Tokenizer Transformer (FT-Transformer) is a simple yet well-performing transformer model for tabular prediction tasks \citep{gorishniy2021revisiting}. The transformer module in FT-Transformer consists of a Multi-Head Self-Attention (MHSA) block and a Feed Forward block \citep{vaswani2017attention}. Recent work has found FT-Transformers to beat other deep learning methods on tabular data \citep{grinsztajn2022tree}.

\textbf{Fastfromer:} Conventional Transformer-like architectures have a quadratic complexity to the length of input sequence \citep{vaswani2017attention}, making them inefficient for tables with large numbers of columns. Fastfromer is an efficient transformer architecture which uses additive attention in place of MHSA \citep{wu2021fastformer}. With additive attention, Fastformer only considers the interaction between each token and the global representation, achieving a linear complexity. 

\textbf{Saint-v:} Saint has introduced the row-wise attention in addition to the column-wise attention of FT-Transformer and Fastformer \citep{somepalli2021saint}. The original implementation of Saint is sensitive to the sequence length and can not handle variable-column tables \citep{somepalli2021saint}. We present a variation of Saint (Saint-v) to fit into our cross-table pretraining setting. Saint-v consists of both column- and row-wise attention blocks, and the detailed model structure is depicted in Appendix \ref{app_saint_v}. 

\subsubsection{Projection heads and objectives} \label{obj} 
There exist various pretraining objectives for tabular prediction tasks \citep{rubachev2022revisiting, majmundar2022met, bahri2021scarf, ucar2021subtab, wang2022transtab, yoon2020vime}. Among them, table reconstruction and contrastive learning are the most popular and effective objectives for tabular tasks. In addition to the self-supervised pretraining objectives, we also tested the pre-finetuning setting using supervised loss.

\textbf{Reconstruction loss:} Reconstruction loss is a self-supervised training objective shown to be effective on various tabular tasks \citep{rubachev2022revisiting, majmundar2022met}. The reconstruction objective aims to recover the original sample $x$ from a corrupted view of the sample $\tilde{x}$. The reconstruction projection head takes the representation of $\tilde{x}$ as input, and generates an estimate of the original input $\hat{x}$. The reconstruction loss is calculated by comparing $x$ and $\hat{x}$. Specifically, we use Cross-Entropy loss to measure the reconstruction error of categorical columns and Mean Squared Error (MSE) for numerical columns.

\textbf{Contrastive loss:} Similar to the reconstruction objective, we also generate $\tilde{x}$ as a corrupted sample. $x$ and its corresponding corruption $\tilde{x}$ are considered as a positive pair of samples, whereas $x$ and other samples in the batch form negative sample pairs. In general, contrastive loss aims to minimize the distance between positive pairs of samples and maximize the distance for negative pairs. Following \citet{bahri2021scarf, chen2020simple}, we used InfoNCE loss for contrastive cross-table pretraining. The contrastive projection heads are similar to those used in SimCLR \citep{chen2020simple}, mapping the representations to the space where we apply the contrastive loss.

\textbf{Supervised loss:} In addition to reconstruction and contrastive losses that do not require labels in pretraining, one can directly pretrain a model using the supervised objective. With supervised losses, the projection head aims to predict the values under a certain field (or column), as predefined by each dataset. The supervised prediction tasks included regression and classification. 

In XTab, the projection heads are data-specific. Different pretraining datasets do not need to share common objectives. For example, we can simultaneously pretrain XTab on both regression and classification tasks, or a mixture of reconstruction and contrastive losses. The diversity of pretraining objectives ensures that the shared backbone is widely adaptable to various downstream tables.

\subsection{Federated pretraining} \label{federated_pretrain} 
XTab introduces data-specific featurizers and projection heads (green blocks in Figure \ref{model_structure}) to account for the variations across table columns and pretraining objectives. During pretraining, both the time and space complexity increase linearly as we include more tabular datasets. As a result, it is challenging to quickly pretrain XTab  using a single machine on a large collection of tabular tasks. To alleviate this issue, we fit XTab into the federated learning framework \citep{mcmahan2017communication}. With the federated setting, XTab involves only marginal overhead in wall-clock time with more pretraining tasks. Federated learning makes it feasible to pretrain XTab on a cluster of commercially available GPUs (NVIDIA T4 GPUs, 16GB memory). 

We use the Federated Averaging (FedAvg) algorithm to pretrain XTab \citep{mcmahan2017communication, li2019convergence}. We have a central server and multiple clients. Each client only hosts one dataset. Therefore, we can distribute the data-specific components of XTab across clients such that each client stores one featurizer, one projection head, and the shared transformer. During pretraining, each client calculates the gradient using the local dataset:

\begin{equation}
\label{eq1}
    w_{k, i+1} \leftarrow w_{k, i}- \alpha \nabla \ell_{k},
\end{equation}

 where $k$ denotes the client (or table) index and $i$ shows the current iteration. $\alpha$ is the learning rate and $\ell^{(k)}$ is the loss function. $w$ represents the trainable parameters which contains two components: $w^{(\text{S})}$ for the shareable modules across all pretraining tasks, and $w^{(\text{NS})}$ for the non-shareable parts ($w  = \text{stack}[w^{(\text{NS})}, w^{(\text{S})}]$). All clients operate synchronously during pretraining with the same learning rate and batch size. 

The central server is responsible for aggregating the local gradients from clients. FedAvg allows clients to make multiple local updates before an aggregation step is made on the central server. Let $N$ denote the number of local updates per aggregation. The central server performs:

\begin{equation}
\label{eq2}
    w^{(\text{S})}_{i+N} \leftarrow w^{(\text{S})}_{i} +  \sum_{k=1}^K (w^{(\text{S})}_{k, i+N} - w^{(\text{S})}_{i}).
\end{equation}

The aggregation is only performed on the shared weights. The term $w^{(\text{S})}_{k, i+N} - w^{(\text{S})}_{i}$ is the gradient learned by client $k$ since the last weight aggregation. The central server simply accumulates the gradients from all clients. Such unitary scalarization was recently shown to perform well in multi-task learning \citep{kurin2022defense}. 

After the aggregation update (i.e., Equation \ref{eq2}), all clients download $w^{(\text{S})}_{i+N}$ from the central server, and apply the weights to the transformer backbone $w_{k, i+N}  = \text{stack}[w_{k, i+N}^{(\text{NS})}, w^{(\text{S})}_{i+N}]$. Therefore, we force all clients to train on a shared backbone with data-specific featurizers and projection heads.

The number of local steps $N$ is a key parameter to control communication efficiency. With $N=1$, FedAvg corresponds to the distributed version of stochastic gradient descent (SGD). With $N>1$, multiple local updates are performed between model aggregation steps at the server, thereby reducing the communication cost between the central server and clients. Unless otherwise specified, we choose $N=5$ throughout the paper. The ablation study on $N$ is shown in Figure \ref{FedAvg} of the Appendix. 

Federated learning was originally proposed as a privacy-preserving approach to learning from distributed data. The collaboration of multiple clients to train a single shared model makes a good fit with our goal of cross-table pretraining. In this work, XTab leverages the distributed nature of federated learning to scale with a large number of pretraining tasks.

\section{Experiments} 
We evaluate the performance of XTab on supervised tabular learning tasks, including binary and multiclass classification and regression. We tested on the following pretraining settings:

\begin{itemize} 
\itemsep-0.3em 
\item XTab with various pretraining objectives, including reconstruction loss, contrastive loss, and supervised loss.
\item XTab with various transformer backbones, including FT-Transformer, Fastformer, and Saint-v. 
\item XTab with the transformer backbone partially- or fully-pretrained from other tasks.
\item XTab with different numbers of pretraining tasks.

\end{itemize}

During finetuning, we randomly initialize a new featurizer and projection head for each downstream task. All downstream tasks use the pretrained transformer backbone. We finetune all the model components using the training set of each downstream task. We included two different finetuning settings:

\begin{itemize} 
\itemsep-0.3em 
\item Light finetuning: finetune XTab for a fixed number of epochs (3 epochs). 
\item Heavy finetuning: finetune XTab with an early stopping patience of 3 epochs. The maximum number of epochs is set to infinity in this case. 

\end{itemize}

For all finetuning settings, we retrieve the best model checkpoint based on validation scores, and use it to report the performance on the test data. The baseline models share the same model architecture and finetuning configurations as XTab, but with randomly initialized parameters instead of using the pretrained backbones. We find that XTab generally outperforms the baseline models in all scenarios and beats other deep learning models on tabular tasks.  Ablation study on the number of pretraining datasets is in Appendix \ref{app_num_tasks}. 

\subsection{Datasets} 
We use the public OpenML-AutoML Benchmark (AMLB: \url{openml.github.io/automlbenchmark/}) \citep{gijsbers2022amlb} for pretraining and evaluation. AMLB is a recently proposed benchmark for automated machine learning, consisting of 104 tabular tasks (71 classification and 33 regression). We included the details of each dataset in Table~\ref{ds_stat} in the Appendix. Out of the 104 tabular datasets, we used 52 datasets for pretraining and the remaining 52 tasks for finetuning and evaluation. We split the pretraining and finetuning datasets by the alphabetical order of the task names (Table~\ref{ds_stat} in the Appendix).  

\textbf{Data split:} For all downstream (or finetuning) tasks, AMLB reserves 10\% of the tabular data for testing. Over the remaining data, we randomly partition 87.5\% (7/8) into the training set and use 12.5\% (1/8) for validation. We repeated 5 trials with different test folds for all tabular datasets. All methods use the same split within the same trial. 

\textbf{Data pre-processing:} Following \citet{bahri2021scarf, somepalli2021saint, wang2022transtab}, we limit the discussion to tables with numerical and categorical columns. Each Category is represented by a distinct integer to index the embedding in the lookup table of the categorical featurizer (see Section \ref{Featurizers} for details). We normalized the numerical features by subtracting the mean and dividing them by the standard deviation. For regression tasks, we also apply the Standardization to the labels. The normalization parameters are calculated using the training set only to avoid information leakage. Missing entries are filled with the mean values of numerical columns, or treated as an additional category for categorical columns. 

\textbf{Table corruption:} Self-supervised learning objectives, including both contrastive and reconstruction losses, require a corrupted view of the input sample. In this work, we follow \citet{bahri2021scarf, rubachev2022revisiting} to randomly resample features and construct a corrupted sample. Specifically, we randomly select a fraction of features at each row of the table. Those features are corrupted by resampling from the empirical marginal distribution of the column. For all datasets, the corruption ratio was set to 60\% as suggested in \citet{bahri2021scarf}. In other words, for each sample $x$ and its corrupted view $\tilde{x}$, 60\% of entries are resampled whereas 40\% of features remain unchanged. 

\subsection{Experimental setup} We used a federated pretraining setting as detailed in Section \ref{federated_pretrain}. Both pretraining and finetuning were performed on a cloud cluster of NVIDIA T4 GPUs (16 GB memory). We used about 30 thousand GPU hours for all experiments.

\textbf{Model configuration and training:} Our default model configuration of transformer variants is the same as \citet{gorishniy2021revisiting}, with 3 transformer blocks, a feature embedding size of 192 and 8 attention heads. The feed forward networks (Figure \ref{model_structure}) have two layers with the same size as the embedding. We apply a dropout ratio of 20\% to attention layers and 10\% for feed forward networks. We use ReGLU \citep{shazeer2020glu} as the activation function and layer normalization \citep{ba2016layer} in the feed forward layers. The projection heads are ReLU networks with 2 layers and a hidden dimension of 192. All model components use \textit{Kaiming} initialization \citep{he2015delving} with the bias terms fixed at zeros. 

The batch size is fixed at 128 for both pretraining and finetuning. Both stages use AdamW as the optimizer, with a learning rate of 1e-4. Following \citet{gorishniy2021revisiting, rubachev2022revisiting}, we also apply a weight decay of 1e-5 to all components excluding featurizers, $\textsf{[CLS]}$ tokens, layer normalization and bias terms.

\textbf{Evaluation metrics:} We choose the evaluation metrics as suggested by AMLB \citep{gijsbers2022amlb}. We use  root mean-squared error (RMSE) for regression tasks, area under the receiver operating characteristic curve (AUC) for binary classification, and log loss for multi-class classification. The same evaluation metrics are applied to validation sets for early stopping. The efficacy of the pretrained transformer backbones is estimated by the downstream performance.

\subsection{Comparison with baseline transformers} 

\textbf{Cross-table pretraining improves downstream task performance.} As shown in Figure \ref{performance}, we compare the downstream prediction performance of FT-Transformer before (baseline) and after cross-table pretraining. Reconstruction objective is used for pretraining and all downstream tasks are finetuned for 3 epochs (light finetuning). We checkpoint the pretrained backbone after a certain number of pretraining steps and finetune downstream tasks from various checkpoints (250/500/1000/1500/2000). In Figure \ref{performance}(a), we show the win rate of the pretrained transformer on all downstream tasks with respect to baseline. Both classification and regression tasks benefit from our proposed cross-table pretraining. As the backbone is pretrained for more steps, we observe an increase in the win rate. We also calculate the rank of the model for each downstream task (Figure \ref{performance}(b)). Model rank is an integer from 1 to 6, with a lower number indicating better performance. Equal values are assigned a rank that is the average of the ranks of those values. The rank of the model improves with XTab pretraining. To further validate the advantage of XTab over transformers without cross-table pretraining, we further look into the normalized prediction performance and error reduction rate (Figure \ref{performance}(c, d)). We min-max normalize the prediction performance of all models, such that the worst model receives a score of 0 and the best model receives 1. Similarly, errors are also normalized to the best and worst models. Negative numbers indicate a model with lower error ($1 - \text{AUC scores}$ for binary classification) or loss (log loss for multiclass classification and RMSE for regression) than baseline. The mean error (or loss) is indicated by the stars. FT-Transformers pretrained with XTab on average obtain higher normalized performance and reduced error compared to traditional random initialization.

\textbf{XTab with different pretraining objectives and finetuning settings.} We extensively test XTab with various pretraining objectives and finetuning settings. Figure \ref{objectives} summarizes the downstream performance using reconstruction, contrastive and supervised objectives as described in Section \ref{obj}. We use FT-Transformer as the backbone. Figure \ref{objectives}(a, b) plot the win rate of XTab under the light and heavy finetuning settings, respectively. We finetune on all downstream tasks for 3 epochs with light finetuning, and use an early stopping patience of 3 for heavy finetuning. We observe a consistent improvement of XTab over the baseline with no cross-table pretraining. The advantage of XTab is more significant in the light finetuning setting compared to heavy finetuning. For example, XTab with the reconstruction objective achieves a 71.0\% win rate with light finetuning, but only 56.1\% with heavy finetuning. The difference is caused by catastrophic forgetting of deep models \citep{ramasesh2021effect, kaushik2021understanding}. As tabular transformers are relatively small ($<$1M parameters for the FT-Transformer backbone), they are more vulnerable to catastrophic forgetting during the finetuning phase. It is possible to alleviate this issue with additional techniques \citep{ramasesh2021effect, kaushik2021understanding}, but this is outside the scope of the paper. Figure \ref{objectives}(c, d) compare different objectives by ranking the models with light and heavy finetuning. All approaches are pretrained for 2000 steps. Each dot in Figure \ref{objectives}(c, d) represents a trial of downstream experiments (5 trials per dataset) and error bars indicate the standard deviations across trials. The advantage of cross-table pretraining is shown by a win rate $>$50\% and a model rank value lower than the baseline. A more detailed comparison involving the normalized performance and error reduction rate is presented in Appendix \ref{app_supp}. We conclude that XTab consistently enhances the downstream performance of tabular transformers across multiple pretraining objectives and finetuning settings. Among all pretraining objectives tested, reconstruction loss performs better than contrastive or supervised losses. 

\textbf{XTab is applicable to various types of transformers.} XTab offers a framework to pretrain the shared model components across tabular tasks. Therefore, the choice of transformer backbone is flexible, as long as the model can process tables with variable columns. In Figure \ref{backbone}, we plug three transformer variants into XTab including FT-Transformer, Fastformer, and Saint-v. The explanation of transformer backbones can be found in Section \ref{method_backbone}. We pretrain all transformers using reconstruction objective, and finetune on the downstream tasks with the light and heavy settings, Figure \ref{backbone}(a, b). We show that XTab is applicable to various types of transformers and all models benefit from the proposed cross-table pretraining, achieving a higher win rate compared to the baseline.

Additional experimental results are presented in the Appendix. In Appendix \ref{appendix_sharable_component}, we pretrain on different components of transformers to identify the shareable components in XTab. In Appendix \ref{app_subsample}, we look into the downstream performance with only a portion of the training set used for finetuning. In Appendix \ref{app_num_tasks}, we compare XTab backbone pretrained on different numbers of tasks and find that more pretraining tasks lead to improved performance. In Appendix \ref{app_fedavg}, we study the federated pretraining setting by changing the number of local updates per global aggregation (i.e., $N$), and find that larger $N$ leads to reduced downstream performance.

\subsection{Performance compared to traditional baselines} 

To compare the performance of XTab and various tabular models, we run experiments on the full AutoML Benchmark \citep{gijsbers2022amlb}. We split the benchmark into 2 folds, each consisting of 52 tabular datasets. We pretrain on fold \#1 and evaluate the downstream performance on fold \#2 and vice versa. We pretrain XTab with the FT-Transformer backbone using reconstruction loss. 20 datasets are excluded since they could not fit into the GPU memory (16 GB, see Table \ref{ds_stat} in the Appendix for details). We report the performance on the remaining 84 tasks. In addition to XTab, we include the following methods:

\textbf{Tree-based models:} Tree-based models provide strong performance on tabular tasks \citep{grinsztajn2022tree}. We include Random Forest (RF) and gradient-boosted tree variants: XGBoost \citep{chen2016xgboost}, LightGBM \citep{ke2017lightgbm} and CatBoost \citep{dorogush2018catboost}. \textbf{Neural networks:} We include the AutoGluon neural networks implemented on top of PyTorch \citep{erickson2020autogluon} and the FastAI tabular model \citep{howard2020fastai}. \textbf{Transformers:} We include the FT-Transformer which is a direct counterpart of XTab without pretraining. The finetuning settings of FTT/XTab include light (FTT-l/XTab-l) and heavy (FTT-h/XTab-h) finetuning as described above. We further introduce FTT-best/XTab-best, which incorporates an early-stopping patience of 20 and model soup of the top 3 checkpoints \citep{wortsman2022model} to achieve better performance. TransTab is included for comparison on classification tasks (regression not enabled yet with TransTab) under the supervised learning (TransTab-sl) and contrastive learning (TransTab-cl) settings \citep{wang2022transtab}. Please refer to Appendix \ref{app_transtab} for how the TransTab ranks are calculated, and Table \ref{t1_class} for results on classification tasks only. 

Table \ref{t1} shows the performance of models with the default hyperparameters and hyperparameter optimization (HPO). With the default hyperparameter, we pretrain XTab for 2000 rounds, whereas the number of pretraining rounds is tuned under the HPO setting. We use the AutoGluon default hyperparameters for tree-based models as they outperform the official defaults to give a strong baseline \citep{erickson2020autogluon}. CatBoost is the state-of-the-art model on tabular tasks, which agrees with the recent finding in \citet{grinsztajn2022tree}. With cross-table pretraining, XTab improves the performance over FTT under light (FTT-l/XTab-l) and heavy (FTT-h/XTab-h) finetuning. Using more finetuning time, XTab-best achieves second place in the benchmark and beats other deep learning models. The success of XTab using the default configuration ensures that the pretrained backbone is widely applicable to tabular tasks, without the need for case-by-case tuning.

With HPO, we randomly search for data-specific hyperparameters on the validation performance. The detailed search space of each model is in Appendix \ref{benchmark_configurations}. We allow a maximum number of 100 HPO trials within a 1-hour time budget. Table \ref{t1} shows that gradient-boosted trees (i.e., XGBoost, LightGBM, CatBoost) achieve higher ranking with HPO, since they are generally faster to train. The search space is also smaller for tree models as they have fewer meaningful hyperparameters and well-known highly performant search spaces. The ranks are calculated separately for default hyperparameters and HPO and  are not comparable across the two settings. The advantage of XTab over FTT increases as we allocate less training time for downstream tasks (XTab-l $\leftarrow$ XTab-h $\leftarrow$ XTab-best $\leftarrow$ XTab with HPO). Therefore, one should use pretrained foundation models instead of randomly initialized weights for tabular transformers, especially with a tight training budget. 

\section{Conclusion} 
In this paper, we present XTab to improve the performance of deep tabular models. XTab pretrains tabular transformers with a diverse collection of data tables, and can improve the tabular prediction performance of an unseen table from arbitrary domains. XTab handles the cross-table variations by separating the models into data-specific and shared components, and encourages the shared components to learn general knowledge for tabular prediction. We also propose to combine self-supervised pretraining with federated learning to improve pretraining efficiency, where client-side nodes perform table reconstruction tasks followed by backbone averaging updates at the server. Our results suggest that finetuning from the pretrained transformer is superior to training tabular transformers from scratch. One limitation of XTab is that it still falls behind CatBoost. This motivates  future works on bridging the gap between pretrained tabular deep learning models and tree models. Another interesting  direction is to combine XTab with language/vision foundation models for improving multimodal learning.

\section*{Software and Data} The AutoML Benchmark (AMLB) is publicly available at \url{openml.github.io/automlbenchmark}. The code and sample pretrained checkpoints are attached to \url{https://github.com/BingzhaoZhu/XTab}.

\end{document}
